% =====================================================================
%  Chapitre 4 — Assistant intelligent et industrialisation
%  Inclusion dans le rapport :  % =====================================================================
%  Chapitre 4 — Assistant intelligent et industrialisation
%  Inclusion dans le rapport :  % =====================================================================
%  Chapitre 4 — Assistant intelligent et industrialisation
%  Inclusion dans le rapport :  % =====================================================================
%  Chapitre 4 — Assistant intelligent et industrialisation
%  Inclusion dans le rapport :  \input{chapitre4}
%
%  Préambule attendu du document principal :
%     \usepackage[T1]{fontenc}  \usepackage[utf8]{inputenc}
%     \usepackage[french]{babel}  \usepackage{graphicx}  \usepackage{booktabs}
%     \graphicspath{{images/}}
%
%  Classe « article » : remplacer \chapter par \section et décaler les
%  \section / \subsection d'un niveau.
%
%  IMAGES — chaque figure porte déjà sa ligne \includegraphics en commentaire.
%  Déposer la capture dans docs/images/ sous le nom indiqué, décommenter cette
%  ligne, puis supprimer le \captureAPlacer situé juste en dessous.
%
%  VOLUME — calibré pour 5 pages en 12 pt, interligne 1,5, avec 6 figures.
%  Si le chapitre déborde : réduire les hauteurs passées à \captureAPlacer
%  (et les largeurs des \includegraphics) de 40mm à 34mm.
% =====================================================================

% Cadre gris affiché tant que la capture n'est pas déposée.
\providecommand{\captureAPlacer}[2]{%
  \fbox{\begin{minipage}[c][#1][c]{0.82\linewidth}%
    \centering\small\itshape Capture à insérer : \texttt{#2}%
  \end{minipage}}%
}

\chapter{Assistant intelligent et industrialisation}

\section{Introduction}

Toutes les réponses de la plateforme décrite jusqu'ici sont calculées par des
requêtes écrites à l'avance. Ce chapitre traite les deux besoins qui échappent
à ce cadre.

Le premier est fonctionnel : « quelles classes ont pris du retard sur le
programme ? » ne désigne aucune colonne, et la circulaire ministérielle
n\textsuperscript{o}~66 du 04/09/2024, qui régit la construction des emplois du
temps, est un document numérisé de sept pages en arabe que personne
n'interroge. Une couche d'assistance conversationnelle a donc été ajoutée, sous
une contrainte explicite : \textbf{aucune réponse n'est émise si elle n'est pas
vérifiable}, et le modèle de langage n'écrit jamais en base. Le second est un
besoin d'industrialisation : six services conteneurisés et trois bases de code
ne se valident pas à la main.

\section{Architecture de la couche intelligente}

L'intelligence artificielle est isolée dans un microservice Python/FastAPI
distinct du backend Java. Ce découpage apporte un écosystème d'outils
(embeddings, extraction PDF) absent de Java, un cycle de vie propre — le
service peut tomber sans emporter la plateforme — et une frontière de sécurité
nette.

Le modèle est \textbf{exécuté localement} par Ollama (\texttt{qwen2.5:7b},
embeddings \texttt{nomic-embed-text}) : aucune donnée d'élève ne quitte
l'établissement. Le choix résulte d'une mesure sur la machine de développement
(RTX 3050, 6~Go) : le 7B tient entièrement en mémoire graphique (4,75~Go) et
répond en 2,7 à 6,5~s, contre 1,4 à 2,7~s pour un 3B nettement moins fiable sur
la traduction de contraintes.

Le service n'a aucun compte de service : il \textbf{rejoue le jeton JWT de
l'utilisateur}, vérifié avec les clés publiques de Keycloak. Un assistant ne
peut donc jamais lire plus que l'utilisateur qui l'interroge. Enfin, le modèle
ne compose jamais lui-même une requête : il \emph{choisit} un outil dans un
catalogue figé de quatorze outils, et le code exécute une requête écrite à la
main. Un mode de défaillance est ainsi supprimé plutôt qu'interdit par une
consigne.

\begin{figure}[htbp]
  \centering
  % \includegraphics[width=0.82\linewidth]{ch4-architecture.png}
  \captureAPlacer{46mm}{images/ch4-architecture.png}
  \caption{Architecture de la couche intelligente : clients web et mobile,
  microservice FastAPI, Ollama local et les trois sources de vérité.}
  \label{fig:ch4-architecture}
\end{figure}

\begin{table}[htbp]
  \centering\small
  \caption{Les trois assistants, leur rôle et leur source de vérité.}
  \label{tab:ch4-assistants}
  \begin{tabular}{@{}llll@{}}
    \toprule
    \textbf{Assistant} & \textbf{Rôle} & \textbf{Source} & \textbf{Écriture} \\
    \midrule
    Contraintes    & Administrateur & API planning + circulaire & sur confirmation \\
    Cahier / cours & Enseignant     & Séances de l'utilisateur, PDF & aucune \\
    Supervision    & Super admin    & Prometheus + Actuator & aucune \\
    \bottomrule
  \end{tabular}
\end{table}

\section{Assistant de l'administrateur : les contraintes}

\subsection{Le problème traité}

La génération d'un emploi du temps repose sur un solveur piloté par des
contraintes exprimées dans un langage dédié (DSL). Ces règles sont le
c\oe{}ur du produit, et pourtant elles restaient inaccessibles au directeur :
les écrire suppose de connaître la syntaxe, le catalogue des types et leurs
paramètres. Les recommandations de la circulaire vivaient donc dans sa tête,
sans jamais devenir des règles activées dans le solveur.

\subsection{La chaîne de traduction}

L'assistant traduit une phrase en règle DSL au terme d'une chaîne dont chaque
étape retire au modèle une décision :

\begin{enumerate}\itemsep1pt \parskip0pt
  \item recherche des trois articles les plus pertinents de la circulaire ;
  \item prompt construit sur le \emph{catalogue réel} du backend et ces articles ;
  \item génération JSON contrainte, température 0 ;
  \item \textbf{validation par le validateur Java}, puis analyse d'impact
        (conflits avec les règles déjà actives) ;
  \item en cas de refus, une seule correction, l'erreur du validateur en entrée ;
  \item restitution de la règle, des conflits et des \textbf{articles cités} —
        \textbf{sans aucune écriture}.
\end{enumerate}

L'enregistrement est une route distincte, déclenchée par un clic sur la règle
affichée mot pour mot. Le modèle ne décide ni de la validité de la règle (le
validateur le fait), ni de sa faisabilité (l'analyse d'impact), ni de son
enregistrement (l'utilisateur). C'est cette séparation qui rend acceptable
qu'un modèle de sept milliards de paramètres se trompe : une règle fausse est
\emph{refusée}, jamais \emph{appliquée}.

\subsection{Ancrage documentaire et évaluation}

Les vingt-deux articles de la circulaire sont indexés une fois pour toutes ;
à cette taille, aucune base vectorielle n'est justifiée. La recherche est
\textbf{hybride}, et c'est une mesure qui l'a imposée : tous les articles
traitant du même sujet dans la même langue administrative, la recherche dense
seule répondait à « combien d'heures de mathématiques en
8\textsuperscript{e} ? » par l'article sur la durée de la journée plutôt que
par le tableau des volumes horaires. Un canal lexical pondéré par l'IDF a été
ajouté, le score valant $\text{dense} + \lambda \cdot \text{lexical}$ avec
$\lambda = 0{,}30$. Sur seize questions de référence, le bon article passe de
\textbf{6 à 11} en première position et de \textbf{10 à 14} dans les trois
premiers.

L'apport dépasse la performance : le directeur qui active « six heures par jour
au maximum » sait que la règle vient de l'article~II.2, page~2, et peut le
lire. La contrainte devient \textbf{traçable jusqu'à sa source réglementaire},
ce qu'aucun validateur ne peut fournir.

\begin{figure}[htbp]
  \centering
  % \includegraphics[width=0.82\linewidth]{ch4-assistant-contraintes.png}
  \captureAPlacer{40mm}{images/ch4-assistant-contraintes.png}
  \caption{Proposition d'une contrainte à partir d'une demande en langage
  naturel : règle DSL, conflits détectés et articles cités.}
  \label{fig:ch4-contraintes}
\end{figure}

\section{Assistant de l'enseignant}

\paragraph{Interroger ses cahiers de séance.} Un cahier de séance est du texte
libre — « les élèves bloquent sur la factorisation ». Aucune requête SQL ne
répond à « qu'ai-je traité avec la 8\textsuperscript{e}~A ? » : la notion
cherchée est dans les mots. Une recherche sémantique est donc construite sur
les séances de l'utilisateur, indexées à la première question puis évincées
après dix minutes. Le même mécanisme sert le directeur pour le pilotage
pédagogique : seul change le corpus, décidé par le backend à partir du compte
porté par le jeton.

\paragraph{Préparer un cours.} L'enseignant joint le chapitre qu'il va traiter
et demande une explication, des exercices ou un QCM corrigé. Le document est
extrait une fois, gardé en mémoire — jamais sur disque — et \textbf{rattaché au
compte qui l'a déposé} : un identifiant deviné n'ouvre pas le cours d'un
collègue. Deux garde-fous complètent l'entrée : un PDF numérisé sans couche de
texte est \emph{refusé en nommant la cause} plutôt que résumé au hasard, et
toute troncature est annoncée à l'utilisateur.

\begin{figure}[htbp]
  \centering
  % \includegraphics[width=0.82\linewidth]{ch4-assistant-enseignant.png}
  \captureAPlacer{40mm}{images/ch4-assistant-enseignant.png}
  \caption{Application mobile : interrogation des cahiers de séance et
  préparation de cours à partir d'un document déposé.}
  \label{fig:ch4-enseignant}
\end{figure}

\section{Assistant du super administrateur}

Le super administrateur exploite une plateforme multi-établissements dont
l'état technique tient dans plusieurs centaines de séries temporelles.
L'assistant de supervision les traduit en français : « y a-t-il un problème sur
la plateforme ? » déclenche des outils tels que \texttt{get\_platform\_health},
\texttt{get\_memory\_usage} ou \texttt{get\_slowest\_endpoints}.

La garantie est la même qu'ailleurs, appliquée ici à la métrologie :
\textbf{les requêtes PromQL sont figées dans le code}, le modèle ne fait que
désigner celle à exécuter. La réponse expose les outils réellement appelés ;
une réponse sans appel d'outil est, par construction, inventée — c'est le
critère de vérification retenu par les tests.

\begin{figure}[htbp]
  \centering
  % \includegraphics[width=0.82\linewidth]{ch4-assistant-supervision.png}
  \captureAPlacer{40mm}{images/ch4-assistant-supervision.png}
  \caption{Assistant de supervision : diagnostic adossé aux métriques réelles,
  avec la liste des outils appelés.}
  \label{fig:ch4-supervision}
\end{figure}

\section{Industrialisation et exploitation}

\subsection{Intégration continue}

Un pipeline GitHub Actions se déclenche à chaque \emph{push} et chaque
\emph{pull request} vers \texttt{main}, en quatre travaux : \textbf{backend}
(\texttt{mvn verify} sur six modules, 845 assertions, couverture JaCoCo agrégée
et analyse SonarCloud) ; \textbf{front} (ESLint puis contrôle de typage
\texttt{tsc}, qui a rattrapé les imports morts lors de la suppression de
modules) ; \textbf{assistant IA} (193 tests \texttt{pytest}, exécutables en
intégration continue parce que Ollama et le backend y sont remplacés par des
doubles) ; \textbf{images Docker} (construction des trois images, sans
publication). Deux réglages méritent d'être signalés : \texttt{concurrency}
annule le pipeline précédent d'une branche encore en cours, et
\texttt{fetch-depth:~0} est indispensable à SonarCloud pour distinguer le code
neuf du code existant.

\subsection{Qualité du code : JaCoCo et SonarCloud}

La couverture est mesurée par JaCoCo module par module, puis agrégée dans un
module dédié. Le rapport est publié en artefact et son résumé — lignes,
branches, instructions — injecté dans le compte rendu de chaque exécution. Ce
même fichier XML est transmis à SonarCloud, qui ajoute les axes que la
couverture ne mesure pas : bogues probables, vulnérabilités, points sensibles
de sécurité, duplications et dette technique.

\begin{figure}[htbp]
  \centering
  % \includegraphics[width=0.82\linewidth]{ch4-ci-sonar.png}
  \captureAPlacer{40mm}{images/ch4-ci-sonar.png}
  \caption{Pipeline d'intégration continue et tableau de bord SonarCloud du
  projet avec la couverture agrégée.}
  \label{fig:ch4-ci}
\end{figure}

\subsection{Conteneurisation et supervision}

Un fichier \texttt{docker-compose.yml} décrit l'environnement complet :
Keycloak et sa base, la base applicative PostgreSQL, Prometheus, Grafana et
l'assistant IA ; une surcouche expose la plateforme à l'application mobile sur
le réseau local. Aucun secret n'est versionné — seul le modèle
\texttt{.env.example} figure au dépôt — et Liquibase reste l'unique source de
vérité du schéma.

Le backend expose ses métriques via Actuator et Micrometer ; Prometheus les
collecte, Grafana les restitue dans un tableau de bord provisionné au
démarrage. Cette chaîne n'est pas décorative : elle est la source de données de
l'assistant de supervision, ce qui fait de l'observabilité une brique
fonctionnelle du produit et non un simple outil d'exploitation.

\begin{figure}[htbp]
  \centering
  % \includegraphics[width=0.82\linewidth]{ch4-monitoring.png}
  \captureAPlacer{40mm}{images/ch4-monitoring.png}
  \caption{Tableau de bord Grafana alimenté par Prometheus : mémoire JVM,
  charge processeur et latence HTTP.}
  \label{fig:ch4-monitoring}
\end{figure}

\section{Conclusion}

Les deux moitiés de ce chapitre servent la même exigence : \emph{une sortie
n'est acceptée que si elle est vérifiable}. Le choix d'outils la vérifie contre
des requêtes écrites à la main, la recherche documentaire contre un texte
citable, la traduction de contraintes contre un validateur Java, et
l'intégration continue vérifie l'ensemble à chaque modification. Les limites
sont assumées : le modèle reste un 7B local, parfois inexact mais toujours
soumis à un contrôle externe avant d'atteindre les données, et la chaîne
s'arrête à la construction des images, le déploiement continu sortant du
périmètre du projet.

%
%  Préambule attendu du document principal :
%     \usepackage[T1]{fontenc}  \usepackage[utf8]{inputenc}
%     \usepackage[french]{babel}  \usepackage{graphicx}  \usepackage{booktabs}
%     \graphicspath{{images/}}
%
%  Classe « article » : remplacer \chapter par \section et décaler les
%  \section / \subsection d'un niveau.
%
%  IMAGES — chaque figure porte déjà sa ligne \includegraphics en commentaire.
%  Déposer la capture dans docs/images/ sous le nom indiqué, décommenter cette
%  ligne, puis supprimer le \captureAPlacer situé juste en dessous.
%
%  VOLUME — calibré pour 5 pages en 12 pt, interligne 1,5, avec 6 figures.
%  Si le chapitre déborde : réduire les hauteurs passées à \captureAPlacer
%  (et les largeurs des \includegraphics) de 40mm à 34mm.
% =====================================================================

% Cadre gris affiché tant que la capture n'est pas déposée.
\providecommand{\captureAPlacer}[2]{%
  \fbox{\begin{minipage}[c][#1][c]{0.82\linewidth}%
    \centering\small\itshape Capture à insérer : \texttt{#2}%
  \end{minipage}}%
}

\chapter{Assistant intelligent et industrialisation}

\section{Introduction}

Toutes les réponses de la plateforme décrite jusqu'ici sont calculées par des
requêtes écrites à l'avance. Ce chapitre traite les deux besoins qui échappent
à ce cadre.

Le premier est fonctionnel : « quelles classes ont pris du retard sur le
programme ? » ne désigne aucune colonne, et la circulaire ministérielle
n\textsuperscript{o}~66 du 04/09/2024, qui régit la construction des emplois du
temps, est un document numérisé de sept pages en arabe que personne
n'interroge. Une couche d'assistance conversationnelle a donc été ajoutée, sous
une contrainte explicite : \textbf{aucune réponse n'est émise si elle n'est pas
vérifiable}, et le modèle de langage n'écrit jamais en base. Le second est un
besoin d'industrialisation : six services conteneurisés et trois bases de code
ne se valident pas à la main.

\section{Architecture de la couche intelligente}

L'intelligence artificielle est isolée dans un microservice Python/FastAPI
distinct du backend Java. Ce découpage apporte un écosystème d'outils
(embeddings, extraction PDF) absent de Java, un cycle de vie propre — le
service peut tomber sans emporter la plateforme — et une frontière de sécurité
nette.

Le modèle est \textbf{exécuté localement} par Ollama (\texttt{qwen2.5:7b},
embeddings \texttt{nomic-embed-text}) : aucune donnée d'élève ne quitte
l'établissement. Le choix résulte d'une mesure sur la machine de développement
(RTX 3050, 6~Go) : le 7B tient entièrement en mémoire graphique (4,75~Go) et
répond en 2,7 à 6,5~s, contre 1,4 à 2,7~s pour un 3B nettement moins fiable sur
la traduction de contraintes.

Le service n'a aucun compte de service : il \textbf{rejoue le jeton JWT de
l'utilisateur}, vérifié avec les clés publiques de Keycloak. Un assistant ne
peut donc jamais lire plus que l'utilisateur qui l'interroge. Enfin, le modèle
ne compose jamais lui-même une requête : il \emph{choisit} un outil dans un
catalogue figé de quatorze outils, et le code exécute une requête écrite à la
main. Un mode de défaillance est ainsi supprimé plutôt qu'interdit par une
consigne.

\begin{figure}[htbp]
  \centering
  % \includegraphics[width=0.82\linewidth]{ch4-architecture.png}
  \captureAPlacer{46mm}{images/ch4-architecture.png}
  \caption{Architecture de la couche intelligente : clients web et mobile,
  microservice FastAPI, Ollama local et les trois sources de vérité.}
  \label{fig:ch4-architecture}
\end{figure}

\begin{table}[htbp]
  \centering\small
  \caption{Les trois assistants, leur rôle et leur source de vérité.}
  \label{tab:ch4-assistants}
  \begin{tabular}{@{}llll@{}}
    \toprule
    \textbf{Assistant} & \textbf{Rôle} & \textbf{Source} & \textbf{Écriture} \\
    \midrule
    Contraintes    & Administrateur & API planning + circulaire & sur confirmation \\
    Cahier / cours & Enseignant     & Séances de l'utilisateur, PDF & aucune \\
    Supervision    & Super admin    & Prometheus + Actuator & aucune \\
    \bottomrule
  \end{tabular}
\end{table}

\section{Assistant de l'administrateur : les contraintes}

\subsection{Le problème traité}

La génération d'un emploi du temps repose sur un solveur piloté par des
contraintes exprimées dans un langage dédié (DSL). Ces règles sont le
c\oe{}ur du produit, et pourtant elles restaient inaccessibles au directeur :
les écrire suppose de connaître la syntaxe, le catalogue des types et leurs
paramètres. Les recommandations de la circulaire vivaient donc dans sa tête,
sans jamais devenir des règles activées dans le solveur.

\subsection{La chaîne de traduction}

L'assistant traduit une phrase en règle DSL au terme d'une chaîne dont chaque
étape retire au modèle une décision :

\begin{enumerate}\itemsep1pt \parskip0pt
  \item recherche des trois articles les plus pertinents de la circulaire ;
  \item prompt construit sur le \emph{catalogue réel} du backend et ces articles ;
  \item génération JSON contrainte, température 0 ;
  \item \textbf{validation par le validateur Java}, puis analyse d'impact
        (conflits avec les règles déjà actives) ;
  \item en cas de refus, une seule correction, l'erreur du validateur en entrée ;
  \item restitution de la règle, des conflits et des \textbf{articles cités} —
        \textbf{sans aucune écriture}.
\end{enumerate}

L'enregistrement est une route distincte, déclenchée par un clic sur la règle
affichée mot pour mot. Le modèle ne décide ni de la validité de la règle (le
validateur le fait), ni de sa faisabilité (l'analyse d'impact), ni de son
enregistrement (l'utilisateur). C'est cette séparation qui rend acceptable
qu'un modèle de sept milliards de paramètres se trompe : une règle fausse est
\emph{refusée}, jamais \emph{appliquée}.

\subsection{Ancrage documentaire et évaluation}

Les vingt-deux articles de la circulaire sont indexés une fois pour toutes ;
à cette taille, aucune base vectorielle n'est justifiée. La recherche est
\textbf{hybride}, et c'est une mesure qui l'a imposée : tous les articles
traitant du même sujet dans la même langue administrative, la recherche dense
seule répondait à « combien d'heures de mathématiques en
8\textsuperscript{e} ? » par l'article sur la durée de la journée plutôt que
par le tableau des volumes horaires. Un canal lexical pondéré par l'IDF a été
ajouté, le score valant $\text{dense} + \lambda \cdot \text{lexical}$ avec
$\lambda = 0{,}30$. Sur seize questions de référence, le bon article passe de
\textbf{6 à 11} en première position et de \textbf{10 à 14} dans les trois
premiers.

L'apport dépasse la performance : le directeur qui active « six heures par jour
au maximum » sait que la règle vient de l'article~II.2, page~2, et peut le
lire. La contrainte devient \textbf{traçable jusqu'à sa source réglementaire},
ce qu'aucun validateur ne peut fournir.

\begin{figure}[htbp]
  \centering
  % \includegraphics[width=0.82\linewidth]{ch4-assistant-contraintes.png}
  \captureAPlacer{40mm}{images/ch4-assistant-contraintes.png}
  \caption{Proposition d'une contrainte à partir d'une demande en langage
  naturel : règle DSL, conflits détectés et articles cités.}
  \label{fig:ch4-contraintes}
\end{figure}

\section{Assistant de l'enseignant}

\paragraph{Interroger ses cahiers de séance.} Un cahier de séance est du texte
libre — « les élèves bloquent sur la factorisation ». Aucune requête SQL ne
répond à « qu'ai-je traité avec la 8\textsuperscript{e}~A ? » : la notion
cherchée est dans les mots. Une recherche sémantique est donc construite sur
les séances de l'utilisateur, indexées à la première question puis évincées
après dix minutes. Le même mécanisme sert le directeur pour le pilotage
pédagogique : seul change le corpus, décidé par le backend à partir du compte
porté par le jeton.

\paragraph{Préparer un cours.} L'enseignant joint le chapitre qu'il va traiter
et demande une explication, des exercices ou un QCM corrigé. Le document est
extrait une fois, gardé en mémoire — jamais sur disque — et \textbf{rattaché au
compte qui l'a déposé} : un identifiant deviné n'ouvre pas le cours d'un
collègue. Deux garde-fous complètent l'entrée : un PDF numérisé sans couche de
texte est \emph{refusé en nommant la cause} plutôt que résumé au hasard, et
toute troncature est annoncée à l'utilisateur.

\begin{figure}[htbp]
  \centering
  % \includegraphics[width=0.82\linewidth]{ch4-assistant-enseignant.png}
  \captureAPlacer{40mm}{images/ch4-assistant-enseignant.png}
  \caption{Application mobile : interrogation des cahiers de séance et
  préparation de cours à partir d'un document déposé.}
  \label{fig:ch4-enseignant}
\end{figure}

\section{Assistant du super administrateur}

Le super administrateur exploite une plateforme multi-établissements dont
l'état technique tient dans plusieurs centaines de séries temporelles.
L'assistant de supervision les traduit en français : « y a-t-il un problème sur
la plateforme ? » déclenche des outils tels que \texttt{get\_platform\_health},
\texttt{get\_memory\_usage} ou \texttt{get\_slowest\_endpoints}.

La garantie est la même qu'ailleurs, appliquée ici à la métrologie :
\textbf{les requêtes PromQL sont figées dans le code}, le modèle ne fait que
désigner celle à exécuter. La réponse expose les outils réellement appelés ;
une réponse sans appel d'outil est, par construction, inventée — c'est le
critère de vérification retenu par les tests.

\begin{figure}[htbp]
  \centering
  % \includegraphics[width=0.82\linewidth]{ch4-assistant-supervision.png}
  \captureAPlacer{40mm}{images/ch4-assistant-supervision.png}
  \caption{Assistant de supervision : diagnostic adossé aux métriques réelles,
  avec la liste des outils appelés.}
  \label{fig:ch4-supervision}
\end{figure}

\section{Industrialisation et exploitation}

\subsection{Intégration continue}

Un pipeline GitHub Actions se déclenche à chaque \emph{push} et chaque
\emph{pull request} vers \texttt{main}, en quatre travaux : \textbf{backend}
(\texttt{mvn verify} sur six modules, 845 assertions, couverture JaCoCo agrégée
et analyse SonarCloud) ; \textbf{front} (ESLint puis contrôle de typage
\texttt{tsc}, qui a rattrapé les imports morts lors de la suppression de
modules) ; \textbf{assistant IA} (193 tests \texttt{pytest}, exécutables en
intégration continue parce que Ollama et le backend y sont remplacés par des
doubles) ; \textbf{images Docker} (construction des trois images, sans
publication). Deux réglages méritent d'être signalés : \texttt{concurrency}
annule le pipeline précédent d'une branche encore en cours, et
\texttt{fetch-depth:~0} est indispensable à SonarCloud pour distinguer le code
neuf du code existant.

\subsection{Qualité du code : JaCoCo et SonarCloud}

La couverture est mesurée par JaCoCo module par module, puis agrégée dans un
module dédié. Le rapport est publié en artefact et son résumé — lignes,
branches, instructions — injecté dans le compte rendu de chaque exécution. Ce
même fichier XML est transmis à SonarCloud, qui ajoute les axes que la
couverture ne mesure pas : bogues probables, vulnérabilités, points sensibles
de sécurité, duplications et dette technique.

\begin{figure}[htbp]
  \centering
  % \includegraphics[width=0.82\linewidth]{ch4-ci-sonar.png}
  \captureAPlacer{40mm}{images/ch4-ci-sonar.png}
  \caption{Pipeline d'intégration continue et tableau de bord SonarCloud du
  projet avec la couverture agrégée.}
  \label{fig:ch4-ci}
\end{figure}

\subsection{Conteneurisation et supervision}

Un fichier \texttt{docker-compose.yml} décrit l'environnement complet :
Keycloak et sa base, la base applicative PostgreSQL, Prometheus, Grafana et
l'assistant IA ; une surcouche expose la plateforme à l'application mobile sur
le réseau local. Aucun secret n'est versionné — seul le modèle
\texttt{.env.example} figure au dépôt — et Liquibase reste l'unique source de
vérité du schéma.

Le backend expose ses métriques via Actuator et Micrometer ; Prometheus les
collecte, Grafana les restitue dans un tableau de bord provisionné au
démarrage. Cette chaîne n'est pas décorative : elle est la source de données de
l'assistant de supervision, ce qui fait de l'observabilité une brique
fonctionnelle du produit et non un simple outil d'exploitation.

\begin{figure}[htbp]
  \centering
  % \includegraphics[width=0.82\linewidth]{ch4-monitoring.png}
  \captureAPlacer{40mm}{images/ch4-monitoring.png}
  \caption{Tableau de bord Grafana alimenté par Prometheus : mémoire JVM,
  charge processeur et latence HTTP.}
  \label{fig:ch4-monitoring}
\end{figure}

\section{Conclusion}

Les deux moitiés de ce chapitre servent la même exigence : \emph{une sortie
n'est acceptée que si elle est vérifiable}. Le choix d'outils la vérifie contre
des requêtes écrites à la main, la recherche documentaire contre un texte
citable, la traduction de contraintes contre un validateur Java, et
l'intégration continue vérifie l'ensemble à chaque modification. Les limites
sont assumées : le modèle reste un 7B local, parfois inexact mais toujours
soumis à un contrôle externe avant d'atteindre les données, et la chaîne
s'arrête à la construction des images, le déploiement continu sortant du
périmètre du projet.

%
%  Préambule attendu du document principal :
%     \usepackage[T1]{fontenc}  \usepackage[utf8]{inputenc}
%     \usepackage[french]{babel}  \usepackage{graphicx}  \usepackage{booktabs}
%     \graphicspath{{images/}}
%
%  Classe « article » : remplacer \chapter par \section et décaler les
%  \section / \subsection d'un niveau.
%
%  IMAGES — chaque figure porte déjà sa ligne \includegraphics en commentaire.
%  Déposer la capture dans docs/images/ sous le nom indiqué, décommenter cette
%  ligne, puis supprimer le \captureAPlacer situé juste en dessous.
%
%  VOLUME — calibré pour 5 pages en 12 pt, interligne 1,5, avec 6 figures.
%  Si le chapitre déborde : réduire les hauteurs passées à \captureAPlacer
%  (et les largeurs des \includegraphics) de 40mm à 34mm.
% =====================================================================

% Cadre gris affiché tant que la capture n'est pas déposée.
\providecommand{\captureAPlacer}[2]{%
  \fbox{\begin{minipage}[c][#1][c]{0.82\linewidth}%
    \centering\small\itshape Capture à insérer : \texttt{#2}%
  \end{minipage}}%
}

\chapter{Assistant intelligent et industrialisation}

\section{Introduction}

Toutes les réponses de la plateforme décrite jusqu'ici sont calculées par des
requêtes écrites à l'avance. Ce chapitre traite les deux besoins qui échappent
à ce cadre.

Le premier est fonctionnel : « quelles classes ont pris du retard sur le
programme ? » ne désigne aucune colonne, et la circulaire ministérielle
n\textsuperscript{o}~66 du 04/09/2024, qui régit la construction des emplois du
temps, est un document numérisé de sept pages en arabe que personne
n'interroge. Une couche d'assistance conversationnelle a donc été ajoutée, sous
une contrainte explicite : \textbf{aucune réponse n'est émise si elle n'est pas
vérifiable}, et le modèle de langage n'écrit jamais en base. Le second est un
besoin d'industrialisation : six services conteneurisés et trois bases de code
ne se valident pas à la main.

\section{Architecture de la couche intelligente}

L'intelligence artificielle est isolée dans un microservice Python/FastAPI
distinct du backend Java. Ce découpage apporte un écosystème d'outils
(embeddings, extraction PDF) absent de Java, un cycle de vie propre — le
service peut tomber sans emporter la plateforme — et une frontière de sécurité
nette.

Le modèle est \textbf{exécuté localement} par Ollama (\texttt{qwen2.5:7b},
embeddings \texttt{nomic-embed-text}) : aucune donnée d'élève ne quitte
l'établissement. Le choix résulte d'une mesure sur la machine de développement
(RTX 3050, 6~Go) : le 7B tient entièrement en mémoire graphique (4,75~Go) et
répond en 2,7 à 6,5~s, contre 1,4 à 2,7~s pour un 3B nettement moins fiable sur
la traduction de contraintes.

Le service n'a aucun compte de service : il \textbf{rejoue le jeton JWT de
l'utilisateur}, vérifié avec les clés publiques de Keycloak. Un assistant ne
peut donc jamais lire plus que l'utilisateur qui l'interroge. Enfin, le modèle
ne compose jamais lui-même une requête : il \emph{choisit} un outil dans un
catalogue figé de quatorze outils, et le code exécute une requête écrite à la
main. Un mode de défaillance est ainsi supprimé plutôt qu'interdit par une
consigne.

\begin{figure}[htbp]
  \centering
  % \includegraphics[width=0.82\linewidth]{ch4-architecture.png}
  \captureAPlacer{46mm}{images/ch4-architecture.png}
  \caption{Architecture de la couche intelligente : clients web et mobile,
  microservice FastAPI, Ollama local et les trois sources de vérité.}
  \label{fig:ch4-architecture}
\end{figure}

\begin{table}[htbp]
  \centering\small
  \caption{Les trois assistants, leur rôle et leur source de vérité.}
  \label{tab:ch4-assistants}
  \begin{tabular}{@{}llll@{}}
    \toprule
    \textbf{Assistant} & \textbf{Rôle} & \textbf{Source} & \textbf{Écriture} \\
    \midrule
    Contraintes    & Administrateur & API planning + circulaire & sur confirmation \\
    Cahier / cours & Enseignant     & Séances de l'utilisateur, PDF & aucune \\
    Supervision    & Super admin    & Prometheus + Actuator & aucune \\
    \bottomrule
  \end{tabular}
\end{table}

\section{Assistant de l'administrateur : les contraintes}

\subsection{Le problème traité}

La génération d'un emploi du temps repose sur un solveur piloté par des
contraintes exprimées dans un langage dédié (DSL). Ces règles sont le
c\oe{}ur du produit, et pourtant elles restaient inaccessibles au directeur :
les écrire suppose de connaître la syntaxe, le catalogue des types et leurs
paramètres. Les recommandations de la circulaire vivaient donc dans sa tête,
sans jamais devenir des règles activées dans le solveur.

\subsection{La chaîne de traduction}

L'assistant traduit une phrase en règle DSL au terme d'une chaîne dont chaque
étape retire au modèle une décision :

\begin{enumerate}\itemsep1pt \parskip0pt
  \item recherche des trois articles les plus pertinents de la circulaire ;
  \item prompt construit sur le \emph{catalogue réel} du backend et ces articles ;
  \item génération JSON contrainte, température 0 ;
  \item \textbf{validation par le validateur Java}, puis analyse d'impact
        (conflits avec les règles déjà actives) ;
  \item en cas de refus, une seule correction, l'erreur du validateur en entrée ;
  \item restitution de la règle, des conflits et des \textbf{articles cités} —
        \textbf{sans aucune écriture}.
\end{enumerate}

L'enregistrement est une route distincte, déclenchée par un clic sur la règle
affichée mot pour mot. Le modèle ne décide ni de la validité de la règle (le
validateur le fait), ni de sa faisabilité (l'analyse d'impact), ni de son
enregistrement (l'utilisateur). C'est cette séparation qui rend acceptable
qu'un modèle de sept milliards de paramètres se trompe : une règle fausse est
\emph{refusée}, jamais \emph{appliquée}.

\subsection{Ancrage documentaire et évaluation}

Les vingt-deux articles de la circulaire sont indexés une fois pour toutes ;
à cette taille, aucune base vectorielle n'est justifiée. La recherche est
\textbf{hybride}, et c'est une mesure qui l'a imposée : tous les articles
traitant du même sujet dans la même langue administrative, la recherche dense
seule répondait à « combien d'heures de mathématiques en
8\textsuperscript{e} ? » par l'article sur la durée de la journée plutôt que
par le tableau des volumes horaires. Un canal lexical pondéré par l'IDF a été
ajouté, le score valant $\text{dense} + \lambda \cdot \text{lexical}$ avec
$\lambda = 0{,}30$. Sur seize questions de référence, le bon article passe de
\textbf{6 à 11} en première position et de \textbf{10 à 14} dans les trois
premiers.

L'apport dépasse la performance : le directeur qui active « six heures par jour
au maximum » sait que la règle vient de l'article~II.2, page~2, et peut le
lire. La contrainte devient \textbf{traçable jusqu'à sa source réglementaire},
ce qu'aucun validateur ne peut fournir.

\begin{figure}[htbp]
  \centering
  % \includegraphics[width=0.82\linewidth]{ch4-assistant-contraintes.png}
  \captureAPlacer{40mm}{images/ch4-assistant-contraintes.png}
  \caption{Proposition d'une contrainte à partir d'une demande en langage
  naturel : règle DSL, conflits détectés et articles cités.}
  \label{fig:ch4-contraintes}
\end{figure}

\section{Assistant de l'enseignant}

\paragraph{Interroger ses cahiers de séance.} Un cahier de séance est du texte
libre — « les élèves bloquent sur la factorisation ». Aucune requête SQL ne
répond à « qu'ai-je traité avec la 8\textsuperscript{e}~A ? » : la notion
cherchée est dans les mots. Une recherche sémantique est donc construite sur
les séances de l'utilisateur, indexées à la première question puis évincées
après dix minutes. Le même mécanisme sert le directeur pour le pilotage
pédagogique : seul change le corpus, décidé par le backend à partir du compte
porté par le jeton.

\paragraph{Préparer un cours.} L'enseignant joint le chapitre qu'il va traiter
et demande une explication, des exercices ou un QCM corrigé. Le document est
extrait une fois, gardé en mémoire — jamais sur disque — et \textbf{rattaché au
compte qui l'a déposé} : un identifiant deviné n'ouvre pas le cours d'un
collègue. Deux garde-fous complètent l'entrée : un PDF numérisé sans couche de
texte est \emph{refusé en nommant la cause} plutôt que résumé au hasard, et
toute troncature est annoncée à l'utilisateur.

\begin{figure}[htbp]
  \centering
  % \includegraphics[width=0.82\linewidth]{ch4-assistant-enseignant.png}
  \captureAPlacer{40mm}{images/ch4-assistant-enseignant.png}
  \caption{Application mobile : interrogation des cahiers de séance et
  préparation de cours à partir d'un document déposé.}
  \label{fig:ch4-enseignant}
\end{figure}

\section{Assistant du super administrateur}

Le super administrateur exploite une plateforme multi-établissements dont
l'état technique tient dans plusieurs centaines de séries temporelles.
L'assistant de supervision les traduit en français : « y a-t-il un problème sur
la plateforme ? » déclenche des outils tels que \texttt{get\_platform\_health},
\texttt{get\_memory\_usage} ou \texttt{get\_slowest\_endpoints}.

La garantie est la même qu'ailleurs, appliquée ici à la métrologie :
\textbf{les requêtes PromQL sont figées dans le code}, le modèle ne fait que
désigner celle à exécuter. La réponse expose les outils réellement appelés ;
une réponse sans appel d'outil est, par construction, inventée — c'est le
critère de vérification retenu par les tests.

\begin{figure}[htbp]
  \centering
  % \includegraphics[width=0.82\linewidth]{ch4-assistant-supervision.png}
  \captureAPlacer{40mm}{images/ch4-assistant-supervision.png}
  \caption{Assistant de supervision : diagnostic adossé aux métriques réelles,
  avec la liste des outils appelés.}
  \label{fig:ch4-supervision}
\end{figure}

\section{Industrialisation et exploitation}

\subsection{Intégration continue}

Un pipeline GitHub Actions se déclenche à chaque \emph{push} et chaque
\emph{pull request} vers \texttt{main}, en quatre travaux : \textbf{backend}
(\texttt{mvn verify} sur six modules, 845 assertions, couverture JaCoCo agrégée
et analyse SonarCloud) ; \textbf{front} (ESLint puis contrôle de typage
\texttt{tsc}, qui a rattrapé les imports morts lors de la suppression de
modules) ; \textbf{assistant IA} (193 tests \texttt{pytest}, exécutables en
intégration continue parce que Ollama et le backend y sont remplacés par des
doubles) ; \textbf{images Docker} (construction des trois images, sans
publication). Deux réglages méritent d'être signalés : \texttt{concurrency}
annule le pipeline précédent d'une branche encore en cours, et
\texttt{fetch-depth:~0} est indispensable à SonarCloud pour distinguer le code
neuf du code existant.

\subsection{Qualité du code : JaCoCo et SonarCloud}

La couverture est mesurée par JaCoCo module par module, puis agrégée dans un
module dédié. Le rapport est publié en artefact et son résumé — lignes,
branches, instructions — injecté dans le compte rendu de chaque exécution. Ce
même fichier XML est transmis à SonarCloud, qui ajoute les axes que la
couverture ne mesure pas : bogues probables, vulnérabilités, points sensibles
de sécurité, duplications et dette technique.

\begin{figure}[htbp]
  \centering
  % \includegraphics[width=0.82\linewidth]{ch4-ci-sonar.png}
  \captureAPlacer{40mm}{images/ch4-ci-sonar.png}
  \caption{Pipeline d'intégration continue et tableau de bord SonarCloud du
  projet avec la couverture agrégée.}
  \label{fig:ch4-ci}
\end{figure}

\subsection{Conteneurisation et supervision}

Un fichier \texttt{docker-compose.yml} décrit l'environnement complet :
Keycloak et sa base, la base applicative PostgreSQL, Prometheus, Grafana et
l'assistant IA ; une surcouche expose la plateforme à l'application mobile sur
le réseau local. Aucun secret n'est versionné — seul le modèle
\texttt{.env.example} figure au dépôt — et Liquibase reste l'unique source de
vérité du schéma.

Le backend expose ses métriques via Actuator et Micrometer ; Prometheus les
collecte, Grafana les restitue dans un tableau de bord provisionné au
démarrage. Cette chaîne n'est pas décorative : elle est la source de données de
l'assistant de supervision, ce qui fait de l'observabilité une brique
fonctionnelle du produit et non un simple outil d'exploitation.

\begin{figure}[htbp]
  \centering
  % \includegraphics[width=0.82\linewidth]{ch4-monitoring.png}
  \captureAPlacer{40mm}{images/ch4-monitoring.png}
  \caption{Tableau de bord Grafana alimenté par Prometheus : mémoire JVM,
  charge processeur et latence HTTP.}
  \label{fig:ch4-monitoring}
\end{figure}

\section{Conclusion}

Les deux moitiés de ce chapitre servent la même exigence : \emph{une sortie
n'est acceptée que si elle est vérifiable}. Le choix d'outils la vérifie contre
des requêtes écrites à la main, la recherche documentaire contre un texte
citable, la traduction de contraintes contre un validateur Java, et
l'intégration continue vérifie l'ensemble à chaque modification. Les limites
sont assumées : le modèle reste un 7B local, parfois inexact mais toujours
soumis à un contrôle externe avant d'atteindre les données, et la chaîne
s'arrête à la construction des images, le déploiement continu sortant du
périmètre du projet.

%
%  Préambule attendu du document principal :
%     \usepackage[T1]{fontenc}  \usepackage[utf8]{inputenc}
%     \usepackage[french]{babel}  \usepackage{graphicx}  \usepackage{booktabs}
%     \graphicspath{{images/}}
%
%  Classe « article » : remplacer \chapter par \section et décaler les
%  \section / \subsection d'un niveau.
%
%  IMAGES — chaque figure porte déjà sa ligne \includegraphics en commentaire.
%  Déposer la capture dans docs/images/ sous le nom indiqué, décommenter cette
%  ligne, puis supprimer le \captureAPlacer situé juste en dessous.
%
%  VOLUME — calibré pour 5 pages en 12 pt, interligne 1,5, avec 6 figures.
%  Si le chapitre déborde : réduire les hauteurs passées à \captureAPlacer
%  (et les largeurs des \includegraphics) de 40mm à 34mm.
% =====================================================================

% Cadre gris affiché tant que la capture n'est pas déposée.
\providecommand{\captureAPlacer}[2]{%
  \fbox{\begin{minipage}[c][#1][c]{0.82\linewidth}%
    \centering\small\itshape Capture à insérer : \texttt{#2}%
  \end{minipage}}%
}

\chapter{Assistant intelligent et industrialisation}

\section{Introduction}

Toutes les réponses de la plateforme décrite jusqu'ici sont calculées par des
requêtes écrites à l'avance. Ce chapitre traite les deux besoins qui échappent
à ce cadre.

Le premier est fonctionnel : « quelles classes ont pris du retard sur le
programme ? » ne désigne aucune colonne, et la circulaire ministérielle
n\textsuperscript{o}~66 du 04/09/2024, qui régit la construction des emplois du
temps, est un document numérisé de sept pages en arabe que personne
n'interroge. Une couche d'assistance conversationnelle a donc été ajoutée, sous
une contrainte explicite : \textbf{aucune réponse n'est émise si elle n'est pas
vérifiable}, et le modèle de langage n'écrit jamais en base. Le second est un
besoin d'industrialisation : six services conteneurisés et trois bases de code
ne se valident pas à la main.

\section{Architecture de la couche intelligente}

L'intelligence artificielle est isolée dans un microservice Python/FastAPI
distinct du backend Java. Ce découpage apporte un écosystème d'outils
(embeddings, extraction PDF) absent de Java, un cycle de vie propre — le
service peut tomber sans emporter la plateforme — et une frontière de sécurité
nette.

Le modèle est \textbf{exécuté localement} par Ollama (\texttt{qwen2.5:7b},
embeddings \texttt{nomic-embed-text}) : aucune donnée d'élève ne quitte
l'établissement. Le choix résulte d'une mesure sur la machine de développement
(RTX 3050, 6~Go) : le 7B tient entièrement en mémoire graphique (4,75~Go) et
répond en 2,7 à 6,5~s, contre 1,4 à 2,7~s pour un 3B nettement moins fiable sur
la traduction de contraintes.

Le service n'a aucun compte de service : il \textbf{rejoue le jeton JWT de
l'utilisateur}, vérifié avec les clés publiques de Keycloak. Un assistant ne
peut donc jamais lire plus que l'utilisateur qui l'interroge. Enfin, le modèle
ne compose jamais lui-même une requête : il \emph{choisit} un outil dans un
catalogue figé de quatorze outils, et le code exécute une requête écrite à la
main. Un mode de défaillance est ainsi supprimé plutôt qu'interdit par une
consigne.

\begin{figure}[htbp]
  \centering
  % \includegraphics[width=0.82\linewidth]{ch4-architecture.png}
  \captureAPlacer{46mm}{images/ch4-architecture.png}
  \caption{Architecture de la couche intelligente : clients web et mobile,
  microservice FastAPI, Ollama local et les trois sources de vérité.}
  \label{fig:ch4-architecture}
\end{figure}

\begin{table}[htbp]
  \centering\small
  \caption{Les trois assistants, leur rôle et leur source de vérité.}
  \label{tab:ch4-assistants}
  \begin{tabular}{@{}llll@{}}
    \toprule
    \textbf{Assistant} & \textbf{Rôle} & \textbf{Source} & \textbf{Écriture} \\
    \midrule
    Contraintes    & Administrateur & API planning + circulaire & sur confirmation \\
    Cahier / cours & Enseignant     & Séances de l'utilisateur, PDF & aucune \\
    Supervision    & Super admin    & Prometheus + Actuator & aucune \\
    \bottomrule
  \end{tabular}
\end{table}

\section{Assistant de l'administrateur : les contraintes}

\subsection{Le problème traité}

La génération d'un emploi du temps repose sur un solveur piloté par des
contraintes exprimées dans un langage dédié (DSL). Ces règles sont le
c\oe{}ur du produit, et pourtant elles restaient inaccessibles au directeur :
les écrire suppose de connaître la syntaxe, le catalogue des types et leurs
paramètres. Les recommandations de la circulaire vivaient donc dans sa tête,
sans jamais devenir des règles activées dans le solveur.

\subsection{La chaîne de traduction}

L'assistant traduit une phrase en règle DSL au terme d'une chaîne dont chaque
étape retire au modèle une décision :

\begin{enumerate}\itemsep1pt \parskip0pt
  \item recherche des trois articles les plus pertinents de la circulaire ;
  \item prompt construit sur le \emph{catalogue réel} du backend et ces articles ;
  \item génération JSON contrainte, température 0 ;
  \item \textbf{validation par le validateur Java}, puis analyse d'impact
        (conflits avec les règles déjà actives) ;
  \item en cas de refus, une seule correction, l'erreur du validateur en entrée ;
  \item restitution de la règle, des conflits et des \textbf{articles cités} —
        \textbf{sans aucune écriture}.
\end{enumerate}

L'enregistrement est une route distincte, déclenchée par un clic sur la règle
affichée mot pour mot. Le modèle ne décide ni de la validité de la règle (le
validateur le fait), ni de sa faisabilité (l'analyse d'impact), ni de son
enregistrement (l'utilisateur). C'est cette séparation qui rend acceptable
qu'un modèle de sept milliards de paramètres se trompe : une règle fausse est
\emph{refusée}, jamais \emph{appliquée}.

\subsection{Ancrage documentaire et évaluation}

Les vingt-deux articles de la circulaire sont indexés une fois pour toutes ;
à cette taille, aucune base vectorielle n'est justifiée. La recherche est
\textbf{hybride}, et c'est une mesure qui l'a imposée : tous les articles
traitant du même sujet dans la même langue administrative, la recherche dense
seule répondait à « combien d'heures de mathématiques en
8\textsuperscript{e} ? » par l'article sur la durée de la journée plutôt que
par le tableau des volumes horaires. Un canal lexical pondéré par l'IDF a été
ajouté, le score valant $\text{dense} + \lambda \cdot \text{lexical}$ avec
$\lambda = 0{,}30$. Sur seize questions de référence, le bon article passe de
\textbf{6 à 11} en première position et de \textbf{10 à 14} dans les trois
premiers.

L'apport dépasse la performance : le directeur qui active « six heures par jour
au maximum » sait que la règle vient de l'article~II.2, page~2, et peut le
lire. La contrainte devient \textbf{traçable jusqu'à sa source réglementaire},
ce qu'aucun validateur ne peut fournir.

\begin{figure}[htbp]
  \centering
  % \includegraphics[width=0.82\linewidth]{ch4-assistant-contraintes.png}
  \captureAPlacer{40mm}{images/ch4-assistant-contraintes.png}
  \caption{Proposition d'une contrainte à partir d'une demande en langage
  naturel : règle DSL, conflits détectés et articles cités.}
  \label{fig:ch4-contraintes}
\end{figure}

\section{Assistant de l'enseignant}

\paragraph{Interroger ses cahiers de séance.} Un cahier de séance est du texte
libre — « les élèves bloquent sur la factorisation ». Aucune requête SQL ne
répond à « qu'ai-je traité avec la 8\textsuperscript{e}~A ? » : la notion
cherchée est dans les mots. Une recherche sémantique est donc construite sur
les séances de l'utilisateur, indexées à la première question puis évincées
après dix minutes. Le même mécanisme sert le directeur pour le pilotage
pédagogique : seul change le corpus, décidé par le backend à partir du compte
porté par le jeton.

\paragraph{Préparer un cours.} L'enseignant joint le chapitre qu'il va traiter
et demande une explication, des exercices ou un QCM corrigé. Le document est
extrait une fois, gardé en mémoire — jamais sur disque — et \textbf{rattaché au
compte qui l'a déposé} : un identifiant deviné n'ouvre pas le cours d'un
collègue. Deux garde-fous complètent l'entrée : un PDF numérisé sans couche de
texte est \emph{refusé en nommant la cause} plutôt que résumé au hasard, et
toute troncature est annoncée à l'utilisateur.

\begin{figure}[htbp]
  \centering
  % \includegraphics[width=0.82\linewidth]{ch4-assistant-enseignant.png}
  \captureAPlacer{40mm}{images/ch4-assistant-enseignant.png}
  \caption{Application mobile : interrogation des cahiers de séance et
  préparation de cours à partir d'un document déposé.}
  \label{fig:ch4-enseignant}
\end{figure}

\section{Assistant du super administrateur}

Le super administrateur exploite une plateforme multi-établissements dont
l'état technique tient dans plusieurs centaines de séries temporelles.
L'assistant de supervision les traduit en français : « y a-t-il un problème sur
la plateforme ? » déclenche des outils tels que \texttt{get\_platform\_health},
\texttt{get\_memory\_usage} ou \texttt{get\_slowest\_endpoints}.

La garantie est la même qu'ailleurs, appliquée ici à la métrologie :
\textbf{les requêtes PromQL sont figées dans le code}, le modèle ne fait que
désigner celle à exécuter. La réponse expose les outils réellement appelés ;
une réponse sans appel d'outil est, par construction, inventée — c'est le
critère de vérification retenu par les tests.

\begin{figure}[htbp]
  \centering
  % \includegraphics[width=0.82\linewidth]{ch4-assistant-supervision.png}
  \captureAPlacer{40mm}{images/ch4-assistant-supervision.png}
  \caption{Assistant de supervision : diagnostic adossé aux métriques réelles,
  avec la liste des outils appelés.}
  \label{fig:ch4-supervision}
\end{figure}

\section{Industrialisation et exploitation}

\subsection{Intégration continue}

Un pipeline GitHub Actions se déclenche à chaque \emph{push} et chaque
\emph{pull request} vers \texttt{main}, en quatre travaux : \textbf{backend}
(\texttt{mvn verify} sur six modules, 845 assertions, couverture JaCoCo agrégée
et analyse SonarCloud) ; \textbf{front} (ESLint puis contrôle de typage
\texttt{tsc}, qui a rattrapé les imports morts lors de la suppression de
modules) ; \textbf{assistant IA} (193 tests \texttt{pytest}, exécutables en
intégration continue parce que Ollama et le backend y sont remplacés par des
doubles) ; \textbf{images Docker} (construction des trois images, sans
publication). Deux réglages méritent d'être signalés : \texttt{concurrency}
annule le pipeline précédent d'une branche encore en cours, et
\texttt{fetch-depth:~0} est indispensable à SonarCloud pour distinguer le code
neuf du code existant.

\subsection{Qualité du code : JaCoCo et SonarCloud}

La couverture est mesurée par JaCoCo module par module, puis agrégée dans un
module dédié. Le rapport est publié en artefact et son résumé — lignes,
branches, instructions — injecté dans le compte rendu de chaque exécution. Ce
même fichier XML est transmis à SonarCloud, qui ajoute les axes que la
couverture ne mesure pas : bogues probables, vulnérabilités, points sensibles
de sécurité, duplications et dette technique.

\begin{figure}[htbp]
  \centering
  % \includegraphics[width=0.82\linewidth]{ch4-ci-sonar.png}
  \captureAPlacer{40mm}{images/ch4-ci-sonar.png}
  \caption{Pipeline d'intégration continue et tableau de bord SonarCloud du
  projet avec la couverture agrégée.}
  \label{fig:ch4-ci}
\end{figure}

\subsection{Conteneurisation et supervision}

Un fichier \texttt{docker-compose.yml} décrit l'environnement complet :
Keycloak et sa base, la base applicative PostgreSQL, Prometheus, Grafana et
l'assistant IA ; une surcouche expose la plateforme à l'application mobile sur
le réseau local. Aucun secret n'est versionné — seul le modèle
\texttt{.env.example} figure au dépôt — et Liquibase reste l'unique source de
vérité du schéma.

Le backend expose ses métriques via Actuator et Micrometer ; Prometheus les
collecte, Grafana les restitue dans un tableau de bord provisionné au
démarrage. Cette chaîne n'est pas décorative : elle est la source de données de
l'assistant de supervision, ce qui fait de l'observabilité une brique
fonctionnelle du produit et non un simple outil d'exploitation.

\begin{figure}[htbp]
  \centering
  % \includegraphics[width=0.82\linewidth]{ch4-monitoring.png}
  \captureAPlacer{40mm}{images/ch4-monitoring.png}
  \caption{Tableau de bord Grafana alimenté par Prometheus : mémoire JVM,
  charge processeur et latence HTTP.}
  \label{fig:ch4-monitoring}
\end{figure}

\section{Conclusion}

Les deux moitiés de ce chapitre servent la même exigence : \emph{une sortie
n'est acceptée que si elle est vérifiable}. Le choix d'outils la vérifie contre
des requêtes écrites à la main, la recherche documentaire contre un texte
citable, la traduction de contraintes contre un validateur Java, et
l'intégration continue vérifie l'ensemble à chaque modification. Les limites
sont assumées : le modèle reste un 7B local, parfois inexact mais toujours
soumis à un contrôle externe avant d'atteindre les données, et la chaîne
s'arrête à la construction des images, le déploiement continu sortant du
périmètre du projet.
